\documentclass[10pt,a4paper]{amsart}
\usepackage{amsmath,amssymb,amsthm,graphicx}
\usepackage{verbatim}
\newenvironment{reference}{\comment}{\endcomment}
%\newenvironment{reference}{}{}
\newenvironment{slogan}{\comment}{\endcomment}
\newenvironment{history}{\comment}{\endcomment}

% For commutative diagrams we use Xy-pic
\usepackage[all]{xy}

% We use 2cell for 2-commutative diagrams.
\xyoption{2cell}
% Theorem environments.
%
\theoremstyle{plain}
\newtheorem{theorem}[subsection]{Theorem}
\newtheorem{proposition}[subsection]{Proposition}
\newtheorem{lemma}[subsection]{Lemma}

\theoremstyle{definition}
\newtheorem{definition}[subsection]{Definition}
\newtheorem{example}[subsection]{Example}
\newtheorem{exercise}[subsection]{Exercise}
\newtheorem{situation}[subsection]{Situation}

\theoremstyle{remark}
\newtheorem{remark}[subsection]{Remark}
\newtheorem{remarks}[subsection]{Remarks}

\numberwithin{equation}{subsection}

% Macros
%
\def\lim{\mathop{\mathrm{lim}}\nolimits}
\def\colim{\mathop{\mathrm{colim}}\nolimits}
\def\Spec{\mathop{\mathrm{Spec}}}
\def\Hom{\mathop{\mathrm{Hom}}\nolimits}
\def\Ext{\mathop{\mathrm{Ext}}\nolimits}
\def\SheafHom{\mathop{\mathcal{H}\!\mathit{om}}\nolimits}
\def\SheafExt{\mathop{\mathcal{E}\!\mathit{xt}}\nolimits}
\def\Sch{\mathit{Sch}}
\def\Mor{\mathop{\mathrm{Mor}}\nolimits}
\def\Ob{\mathop{\mathrm{Ob}}\nolimits}
\def\Sh{\mathop{\mathit{Sh}}\nolimits}
\def\NL{\mathop{N\!L}\nolimits}
\def\CH{\mathop{\mathrm{CH}}\nolimits}
\def\proetale{{pro\text{-}\acute{e}tale}}
\def\etale{{\acute{e}tale}}
\def\QCoh{\mathit{QCoh}}
\def\Ker{\mathop{\mathrm{Ker}}}
\def\Im{\mathop{\mathrm{Im}}}
\def\Coker{\mathop{\mathrm{Coker}}}
\def\Coim{\mathop{\mathrm{Coim}}}

% Boxtimes
%
\DeclareMathSymbol{\boxtimes}{\mathbin}{AMSa}{"02}

%
% Macros for moduli stacks/spaces
%
\def\QCohstack{\mathcal{QC}\!\mathit{oh}}
\def\Cohstack{\mathcal{C}\!\mathit{oh}}
\def\Spacesstack{\mathcal{S}\!\mathit{paces}}
\def\Quotfunctor{\mathrm{Quot}}
\def\Hilbfunctor{\mathrm{Hilb}}
\def\Curvesstack{\mathcal{C}\!\mathit{urves}}
\def\Polarizedstack{\mathcal{P}\!\mathit{olarized}}
\def\Complexesstack{\mathcal{C}\!\mathit{omplexes}}
% \Pic is the operator that assigns to X its picard group, usage \Pic(X)
% \Picardstack_{X/B} denotes the Picard stack of X over B
% \Picardfunctor_{X/B} denotes the Picard functor of X over B
\def\Pic{\mathop{\mathrm{Pic}}\nolimits}
\def\Picardstack{\mathcal{P}\!\mathit{ic}}
\def\Picardfunctor{\mathrm{Pic}}
\def\Deformationcategory{\mathcal{D}\!\mathit{ef}}
\usepackage{cleveref}
\crefname{originalequation}{original equation}{original equations}
\Crefname{originalequation}{Original equation}{Original equations}
\makeatletter
\newlabel{equation-C-prime}{{0AL4}{source}{}{sourcepointer-1}{}}
\newlabel{equation-C-prime@cref}{{[originalequation][0][]0AL4}{[1][0][]source}{}{}{}}
\makeatother
\crefname{originalreference}{original reference}{original references}
\Crefname{originalreference}{Original reference}{Original references}
\makeatletter
\newlabel{lemma-algebraize-easy}{{0GAS}{source}{}{sourcepointer-2}{}}
\newlabel{lemma-algebraize-easy@cref}{{[originalreference][0][]0GAS}{[1][0][]source}{}{}{}}
\makeatother
\crefname{originalreference}{original reference}{original references}
\Crefname{originalreference}{Original reference}{Original references}
\makeatletter
\newlabel{lemma-close-enough}{{0GAR}{source}{}{sourcepointer-3}{}}
\newlabel{lemma-close-enough@cref}{{[originalreference][0][]0GAR}{[1][0][]source}{}{}{}}
\makeatother
\crefname{originallemma}{original lemma}{original lemmas}
\Crefname{originallemma}{Original lemma}{Original lemmas}
\makeatletter
\newlabel{more-algebra-lemma-henselization}{{0A02}{source}{}{sourcepointer-4}{}}
\newlabel{more-algebra-lemma-henselization@cref}{{[originallemma][0][]0A02}{[1][0][]source}{}{}{}}
\makeatother
\crefname{originallemma}{original lemma}{original lemmas}
\Crefname{originallemma}{Original lemma}{Original lemmas}
\makeatletter
\newlabel{more-algebra-lemma-henselization-Noetherian-pair}{{0AGV}{source}{}{sourcepointer-5}{}}
\newlabel{more-algebra-lemma-henselization-Noetherian-pair@cref}{{[originallemma][0][]0AGV}{[1][0][]source}{}{}{}}
\makeatother
\crefname{originalproposition}{original proposition}{original propositions}
\Crefname{originalproposition}{Original proposition}{Original propositions}
\makeatletter
\newlabel{more-algebra-proposition-finite-type-over-G-ring}{{07PV}{source}{}{sourcepointer-6}{}}
\newlabel{more-algebra-proposition-finite-type-over-G-ring@cref}{{[originalproposition][0][]07PV}{[1][0][]source}{}{}{}}
\makeatother
\crefname{originallemma}{original lemma}{original lemmas}
\Crefname{originallemma}{Original lemma}{Original lemmas}
\makeatletter
\newlabel{smoothing-lemma-henselian-pair}{{0AH5}{source}{}{sourcepointer-7}{}}
\newlabel{smoothing-lemma-henselian-pair@cref}{{[originallemma][0][]0AH5}{[1][0][]source}{}{}{}}
\makeatother
\title{Rig-smooth complete algebras over a G-ring are completed finite type algebras}
\author{The Stacks Project Authors}
\date{stacks-project:restricted.tex:local-snapshot-20261012; GNU FDL 1.2 or later}
\begin{document}
\maketitle
\noindent\textit{Editorial note.} All delivered mathematical text is the author's own, copied from the pinned Stacks Project chapter file restricted.tex; no statement or proof step was rewritten, reordered or removed.\par\medskip
\noindent\textit{Editorial note.} The counted proof is the author's complete argument: the proof is complete, from the free resolution of the presentation to the final isomorphism via the close-enough threshold\par\medskip
\noindent\textit{Editorial note.} The label proposition-approximate names the counted statement, and the proof environment immediately following it in the source is delivered verbatim as its proof.\par\medskip
\noindent\textit{Editorial note.} The counted statement is carried by the author's own proposition environment in the chapter source.\par\medskip
\noindent\textit{Allowed standard prerequisites (editorial).}\begin{itemize}
\item presentations of complete algebras and free resolutions of the quotient
\item the henselization of a Noetherian pair and its G-ring stability under polynomial rings
\item Artin approximation for the henselization to completion map of a G-ring pair
\item the close-enough lemma giving a congruence threshold under which quotient completions are isomorphic
\item the easy algebraization lemma identifying a completion of a quotient of the henselization with a completed finite type A-algebra
\end{itemize}\medskip
\ifdefined\hypertarget\hypertarget{sourcepointer-1}{}\fi
\noindent\textit{Source reference.} Original equation 0AL4: sources/stacks-restricted/pinned-original-source.tex, source lines 326--326. Prerequisite pointer invoked at proof line(s) 1477; it names a standard fact of the same chapter that the counted argument uses, and every step of the counted proof built on it is retained verbatim in the delivered body.\par\medskip
\ifdefined\hypertarget\hypertarget{sourcepointer-2}{}\fi
\noindent\textit{Source reference.} Original reference 0GAS: sources/stacks-restricted/pinned-original-source.tex, source lines 1452--1452. Prerequisite pointer invoked at proof line(s) 1528; it names a standard fact of the same chapter that the counted argument uses, and every step of the counted proof built on it is retained verbatim in the delivered body.\par\medskip
\ifdefined\hypertarget\hypertarget{sourcepointer-3}{}\fi
\noindent\textit{Source reference.} Original reference 0GAR: sources/stacks-restricted/pinned-original-source.tex, source lines 1363--1363. Prerequisite pointer invoked at proof line(s) 1515; it names a standard fact of the same chapter that the counted argument uses, and every step of the counted proof built on it is retained verbatim in the delivered body.\par\medskip
\ifdefined\hypertarget\hypertarget{sourcepointer-4}{}\fi
\noindent\textit{Source reference.} Original lemma 0A02: sources/stacks-chow/tags.txt, source lines 12620--12620. Prerequisite pointer invoked at proof line(s) 1503; it names a standard fact of a sibling chapter that the counted argument uses, and every step of the counted proof built on it is retained verbatim in the delivered body.\par\medskip
\ifdefined\hypertarget\hypertarget{sourcepointer-5}{}\fi
\noindent\textit{Source reference.} Original lemma 0AGV: sources/stacks-chow/tags.txt, source lines 13208--13208. Prerequisite pointer invoked at proof line(s) 1507; it names a standard fact of a sibling chapter that the counted argument uses, and every step of the counted proof built on it is retained verbatim in the delivered body.\par\medskip
\ifdefined\hypertarget\hypertarget{sourcepointer-6}{}\fi
\noindent\textit{Source reference.} Original proposition 07PV: sources/stacks-chow/tags.txt, source lines 9810--9810. Prerequisite pointer invoked at proof line(s) 1510; it names a standard fact of a sibling chapter that the counted argument uses, and every step of the counted proof built on it is retained verbatim in the delivered body.\par\medskip
\ifdefined\hypertarget\hypertarget{sourcepointer-7}{}\fi
\noindent\textit{Source reference.} Original lemma 0AH5: sources/stacks-chow/tags.txt, source lines 13218--13218. Prerequisite pointer invoked at proof line(s) 1514; it names a standard fact of a sibling chapter that the counted argument uses, and every step of the counted proof built on it is retained verbatim in the delivered body.\par\medskip
\begin{proposition}
\label{proposition-approximate}
Let $I$ be an ideal of a Noetherian G-ring $A$. Let $B$ be an
object of (\ref{equation-C-prime}). If $B$ is rig-smooth
over $(A, I)$, then there exists a finite type $A$-algebra
$C$ and an isomorphism $B \cong C^\wedge$ of $A$-algebras.
\end{proposition}

\begin{proof}
Choose a presentation $B = A[x_1, \ldots, x_r]^\wedge/J$. Write
$P = A[x_1, \ldots, x_r]^\wedge$. Choose generators $g_1, \ldots, g_m \in J$.
Choose generators $k_1, \ldots, k_t$ of the module
of relations between $g_1, \ldots, g_m$, i.e., such that
$$
P^{\oplus t} \xrightarrow{k_1, \ldots, k_t}
P^{\oplus m} \xrightarrow{g_1, \ldots, g_m}
P \to B \to 0
$$
is a resolution. Write $k_i = (k_{i1}, \ldots, k_{im})$ so that we have
\begin{equation}
\label{equation-relations-straight-up}
\sum\nolimits_j k_{ij}g_j = 0
\end{equation}
for $i = 1, \ldots, t$. Denote $K = (k_{ij})$ the $m \times t$-matrix
with entries $k_{ij}$.

\medskip\noindent
Let $A[x_1, \ldots, x_r]^h$ be the henselization of the
pair $(A[x_1, \ldots, x_r], IA[x_1, \ldots, x_r])$, see
More on Algebra, Lemma \ref{more-algebra-lemma-henselization}.
We may and do think of $A[x_1, \ldots, x_r]^h$ as a subring
of $P = A[x_1, \ldots, x_r]^\wedge$, see
More on Algebra, Lemma
\ref{more-algebra-lemma-henselization-Noetherian-pair}.
Since $A$ is a Noetherian G-ring, so is $A[x_1, \ldots, x_r]$, see
More on Algebra, Proposition
\ref{more-algebra-proposition-finite-type-over-G-ring}.
Hence we have approximation for the map
$A[x_1, \ldots, x_r]^h \to A[x_1, \ldots, x_r]^\wedge = P$
with respect to the ideal generated by $I$, see
Smoothing Ring Maps, Lemma \ref{smoothing-lemma-henselian-pair}.
Choose a large enough integer $n$ as in Lemma \ref{lemma-close-enough}.
By the approximation property we may choose
$g'_1, \ldots, g'_m \in A[x_1, \ldots, x_r]^h$
and a matrix
$K' = (k'_{ij}) \in \text{Mat}(m \times t, A[x_1, \ldots, x_r]^h)$
such that $\sum\nolimits_j k'_{ij}g'_j = 0$ in
$A[x_1, \ldots, x_r]^h$ and such that
$g_i - g'_i \in I^nP$ and $K - K' \in I^n\text{Mat}(m \times t, P)$.
By our choice of $n$ we conclude that there is an isomorphism
$$
B \to P/(g'_1, \ldots, g'_m) =
\left(A[x_1, \ldots, x_r]^h/(g'_1, \ldots, g'_m)\right)^\wedge
$$
This finishes the proof by Lemma \ref{lemma-algebraize-easy}.
\end{proof}

\end{document}
